\section{Exercise Solutions}

\begin{solution}
The AGM iteration $a_{n+1} = (a_n+b_n)/2$, $b_{n+1} = \sqrt{a_n b_n}$ converges quadratically to $M(a_0, b_0)$. For $a_0 = 1, b_0 = 1/\sqrt{2}$, Gauss discovered $M(1, 1/\sqrt{2}) = \frac{\pi}{2\int_0^{\pi/2}\frac{d\theta}{\sqrt{1-\frac{1}{2}\sin^2\theta}}}$, relating the AGM to elliptic integrals. This gives a fast algorithm for computing $\pi$.
\end{solution}

\begin{solution}
The Borwein brothers showed that the AGM can be used to compute $\pi$ with quartic convergence: starting from $a_0 = 1, b_0 = 0, s_0 = 1/4$, the iteration $a_{n+1} = \frac{a_n+\sqrt{a_n b_n}}{2}$, $b_{n+1} = \sqrt{a_n b_n}$, $s_{n+1} = s_n - 2^n(a_n - a_{n+1})^2$ gives $p_n = \frac{4a_n^2}{1-s_n}$ converging to $\pi$ with 4x more digits per iteration.
\end{solution}

\begin{solution}
The AGM relates to elliptic integrals via $\int_0^{\pi/2}\frac{d\theta}{\sqrt{a^2\cos^2\theta + b^2\sin^2\theta}} = \frac{\pi}{2M(a,b)}$. Setting $a = 1, b = k'$ gives the complete elliptic integral $K(k) = \frac{\pi}{2M(1,k')}$ where $k' = \sqrt{1-k^2}$. This was Gauss's key insight connecting his AGM work to elliptic functions.
\end{solution}

\begin{solution}
The descending Landen transformation $k_1 = \frac{1-k'}{1+k'}$ corresponds to one AGM iteration. Iterating this gives a sequence $k_n \to 1$ with $K(k_n) = \frac{1+k_n'}{2}K(k_{n+1})$. The product telescopes to give $K(k) = \frac{\pi}{2M(1,k')}$, proving the AGM-elliptic integral relation.
\end{solution}
